% 148-13(148쪽) — 평행사변형 ABCD · $\seg{AB}=6$, $\seg{BC}=8$.
% 스캔 화질이 낮아 TikZ 로 다시 그림. 원문과 같이 B 를 왼쪽 아래, C 를 오른쪽 아래, A·D 를 그 위에 두었다.
% 라벨은 원문에 있는 A·B·C·D·$6$·$8$ 만 둔다 — 넓이와 $D$ 의 크기(답)는 적지 않는다.
\begin{tikzpicture}[line width=0.5pt, scale=0.68, >=stealth]
  \coordinate (B) at (0,0);
  \coordinate (C) at (3.4,0);
  \coordinate (A) at (0.95,1.6);
  \coordinate (D) at (4.35,1.6);
  \draw[line width=0.55pt] (A) -- (B) -- (C) -- (D) -- cycle;
  \draw[densely dotted, line width=0.35pt, <->] ([shift=(149:0.18)]B) to[bend left=12]
    node[midway, left=-1.5pt, font=\footnotesize] {$6$} ([shift=(149:0.18)]A);
  \draw[densely dotted, line width=0.35pt, <->] ([shift=(-90:0.18)]B) to[bend right=12]
    node[midway, below=-2.5pt, font=\footnotesize] {$8$} ([shift=(-90:0.18)]C);
  \node[above left=-3.5pt and -3pt, font=\small] at (A) {$\pt{A}$};
  \node[left=1pt, font=\small] at (B) {$\pt{B}$};
  \node[below right=-3pt and -3pt, font=\small] at (C) {$\pt{C}$};
  \node[above right=-3.5pt and -3pt, font=\small] at (D) {$\pt{D}$};
\end{tikzpicture}
